\documentclass[12pt]{article}
\usepackage{color,soul}
% Import preambles and macros for notes
\input{../../../../preambles/notes-preamble.tex}
\input{../../../../preambles/notes-macros.tex}
\input{../../../../preambles/theorem-system-section.tex}

\title{MATH 481 Lecture 8}
\author{Deepak Jassal}
\date{February 3\textsuperscript{rd}, 2026}

\begin{document}
\setcounter{section}{3}
\setcounter{tcb@cnt@mydefinition}{1}
\maketitle
\pf[Continue Proof of Theorem 3.1]{
    \begin{enumerate}
        \item[(1)] We are working to prove that $\exists c_1,c_2>0$ such that 
        \[
            c_1x\leq\Psi(x)\leq c_2x
        \]
        for all $x\geq2$.\\
        Last time we showed that given any $\varepsilon>0$, there exists $x_0\geq1$ (possibly $x_0=x_0(\varepsilon)$) such that 
        \[
            (\log2-\varepsilon)x\leq D(x)\leq(\log2+\varepsilon)x
        \]
        for all $x\geq x_0$. We also showed that 
        \[
            \Psi(x)-\Psi\left(\frac{x}{2}\right)\leq D(x)\leq\Psi(x)
        \]
        for $x\geq2$. This gives us 
        \[
            |Psi(x)\geq D(x)\geq(\log2-\varepsilon)x
        \]
        for all $x\geq x_0$. We also have that 
        \begin{align*}
            \Psi(x)&\leq D(x)+\Psi\left(\frac{x}{2}\right)\\
            &\leq(\log 2-\varepsilon)x+\Psi\left(\frac{x}{2}\right)& \forall x\geq x_0
            \intertext{By choosing a small enough $\varepsilon$}
            &\leq x+\Psi\left(\frac{x}{2}\right)& \forall x\geq x_0\\
            &\leq x+\frac{x}{2}+\Psi\left(\frac{x}{4}\right)& \forall \frac{x}{2}\geq x_0\\
            &\quad\quad\quad\quad\:\,\vdots\\
            &\leq x+\frac{x}{2}+\frac{x}{4}+\cdots+\Psi\left(\frac{x}{2^k}\right)& \forall \frac{x}{2^{k-1}}\geq x_0.
        \end{align*}
        Choose $k$ to be the largest integer such that 
        \[
            \frac{x}{2^{k-1}}\geq x_0
        \]
        i.e.,
        \[
            \frac{x}{2^k}\leq x_0\leq\frac{x}{2^{k-1}}.
        \]
        Since $\Psi$ is an increasing function, we have 
        \[
            \Psi(x_0)\geq\Psi\left(\frac{x}{2}\right).
        \]
        It follows that 
        \begin{align*}
            \Psi(x)&\leq x+\frac{x}{2}+\cdots+\frac{x}{2^{k-1}}+\Psi(x_0)\\
            &=x\sum_{i=0}^{k-1}\left(\frac{1}{2}\right)^i+\Psi(x_0)\\
            &=2x+\Psi(x_0)
            \intertext{For large enough $x$ we have $\Psi(x_0)\leq \varepsilon x$}
            \Psi(x)&\leq(2+\varepsilon)x
            \intertext{For all $x\geq X_0$, and some $X_0>1$.}
        \end{align*}
        Combining everything, we get
        \[
            (2-\varepsilon)\leq\Psi(x)\leq(2+\varepsilon)x
        \]
        for all $x\geq X_0$.\\
        What if $2\leq x\leq X_0$? Trivially we have that 
        \[
            \frac{\Psi(x)}{x}\leq\frac{\Psi(X_0)}{2}.
        \]
        We also have that 
        \[
            \frac{\Psi(x)}{x}\geq\frac{\Psi(2)}{x}=\frac{\log2}{x}\leq\frac{\log2}{X_0}.
        \]
        This gives
        \[
            \frac{\log2}{X_0}x\leq\Psi(x)\leq\frac{\Psi(X_0)}{2}x
        \]
        for all $2\leq x\leq X_0$. All in all, $\exists c_1,c_2>0$ such that 
        \[
            c_1x\leq \Psi(x)\leq c_2x
        \]
        for all $x\geq2$.
        \item[(2)] 
        \begin{align*}
            \Psi(x)-\theta(x)&=\sum_{n\leq x}\Lambda(n)-\sum_{p\leq x}\log p\\
            &=\sum_{\substack{n\leq x\\ n=p^m}}\log p-\sum_{p\leq x}\log p\\    
            &=\sum_{\substack{p^m\leq x\\ p\text{ prime}\\m\in\N}}\log p-\sum_{p\leq x}\log p\\
            &=\sum_{\substack{p^m\leq x\\ p\text{ prime}\\m\geq2}}\log p\\
            &=\sum_{p\leq\sqrt{x}}\sum_{\substack{m\leq\frac{\log x}{\log p}\\m\geq2}}\log p\\    
            &=\sum_{p\leq\sqrt{x}}\log p\floor{\frac{\log x}{\log p}}\\    
            &\leq\sum_{p\leq\sqrt{x}}\log p\frac{\log x}{\log p}\\
            &=\log x\sum_{p\leq\sqrt{x}}1\\
            &\leq(\log x)\sqrt{x}.
        \end{align*}
        We just showed that
        \[
            \Psi(x)-\theta(x)\leq (\log x)\sqrt{x}.
        \]
        No we have
        \[
            \theta(x)\leq \Psi(x)\leq c_2x
        \]
        we also have
        \begin{align*}
            \theta(x)&\geq\Psi(x)-(\log x)\sqrt{x}\\
            &\geq c_1x-(\log x)\sqrt{x}
        \end{align*}
        for $x$ large enough we have
        \[
            (\log x)\sqrt{x}\leq\frac{c_1}{2}x
        \]
        and so 
        \[
            \theta(x)\geq c_1x-\frac{c_1}{2}x=\frac{c_1}{2}x.
        \]
        \item[(3)] 
        \begin{align*}
            \pi(x)&=\sum_{p\leq x}1\\
            &\geq\sum_{p\leq x}\frac{\log p}{\log x}\\
            &=\frac{1}{\log x}\sum_{p\leq x}\log p\\
            &=\frac{1}{\log x}\theta(x)\geq\frac{\frac{c_1}{x}x}{\log x}\\
            &=\sum_{p\leq\sqrt{x}}1+\sum_{\sqrt{x}<p\leq x}1\\
            &\leq\sqrt{x}+\frac{1}{\log\sqrt{x}}\sum_{\sqrt{x}<p\leq x}\log\sqrt{x}\\
            &\sqrt{x}+\frac{2}{\log x}\sum_{\sqrt{}\leq p\leq x}\log p\\
            &\leq\sqrt{x}+\frac{2}{\log x}(\theta(x)-\theta(\sqrt{x}))\\
            &\leq\sqrt{x}+\frac{2}{\log x}(c_2x-\frac{c_1}{2}\sqrt{x})\\
            &\leq C\frac{x}{\log x}
        \end{align*}
        for some $C>0$. Indeed we have that 
        \[
            \pi(x)\asymp\frac{x}{\log x}
        \]
        for $x$ large enough. Repeating the argument we did to bound $\Psi(x)$ for small values of $x$, we get $\pi(x)\asymp\frac{x}{\log x}$ for all $x\geq 2$. 
    \end{enumerate}
    Thus, we have all of the desired results.
}
\thm{}{
    For all $x\leq2$, we have 
    \begin{enumerate}
        \item[(1)] $\theta(x)=\Psi(x)+\Oh(\sqrt{x})$;
        \item[(2)] $\pi(x)=\frac{\Psi(x)}{\log x}+\Oh\left(\frac{x}{\log^2x}\right)$.
    \end{enumerate}
}
\pf[Proof of Theorem 3.2]{
    \begin{enumerate}
        \item[(1)]
        \begin{align*}
            \Psi(x)-\theta(x)&=\vdots\\
            &\leq\log x\sum_{p\leq\sqrt{x}}1\\
            &=\log x\pi(\sqrt{x})
            \intertext{By theorem 3.1 part 3}
            &\leq\log x\frac{c_2\sqrt{x}}{\log\sqrt{x}}\\
            &\leq2c_2\sqrt{x}\\
            &\ll\sqrt{x}.
        \end{align*}
        \item[(2)]
        \begin{align*}
            \pi(x)&=\sum_{p\leq x}1\\
            &=\sum_{p\leq x}\frac{\log p}{\log p}\\
            &=\left(\sum_{p\leq x}\log p\right)\frac{1}{\log x}+\int_{2}^{x}\left(\sum_{p\leq t}\log p\right)\frac{1}{t(\log t)^2}\,dt\\
            &=\frac{\Oh(x)}{\log x}+\int_{2}^{x}\frac{\Oh(t)}{t(\log t)^2}\,dt.
        \end{align*}
        Observe that
        \[
            \int_{2}^{x}\frac{\Oh(t)}{t(\log t)^2}\,dt\ll\int_{2}^{x}\frac{1}{(\log t)^2}\,dt=\Li_2(x)\ll\frac{x}{(\log x)^2}.
        \]
        Hence,
        \begin{align*}
            \pi(x)&=\frac{\Oh(x)}{\log x}+\Oh\left(\frac{x}{(\log x)^2}\right)\\
            &=\frac{\Psi(x)}{\log x}+\Oh\left(\frac{\sqrt{x}}{\log x}\right)+\Oh\left(\frac{x}{(\log x)^2}\right)\\
            &=\frac{\Psi(x)}{\log x}+\Oh\left(\frac{x}{(\log x)^2}\right).
        \end{align*}
    \end{enumerate}
    Thus, we have out desired results.
}
\cor{
    The following are all equivalent as $x\to\infty$
    \begin{enumerate}
        \item[(1)] $\pi(x)\sim\frac{x}{\log x}$;
        \item[(2)] $\theta(x)\sim x$;
        \item[(3)] $\Psi(x)\sim x$.   
    \end{enumerate}
}
\section*{Mertens' Estimate}
\thm{}{
    We have
    \begin{enumerate}
        \item[(1)] 
        \[
            \sum_{n\leq x}\frac{\Lambda(n)}{n}=\log x+\Oh(1)
        \]
        \item[(2)]
        \[
            \sum_{p\leq x}\frac{\log p}{p}=\log x+\Oh(1)
        \]
        \item[(3)]
        \[
            \sum_{p\leq x}\frac{1}{p}=\log\log x+A+\Oh\left(\frac{1}{\log x}\right)
        \]
        \item[(4)] Mertens' formula
        \[
            \prod_{p\leq x}\left(1-\frac{1}{p}\right)=\frac{e^{-\gamma}}{\log x}\left(1+\Oh\left(\frac{1}{\log x}\right)\right).
        \]
    \end{enumerate}
}
\pf[Proof of Theorem 3.4]{
    \begin{enumerate}
        \item[(1)] Let 
        \begin{align*}
            S(x)&=\sum_{n\leq x}\log n\\
            &=x\log x-x+\Oh(\log x),
        \end{align*}
        on the other hand we have
        \begin{align*}
            S(x)&=\sum_{d\leq x}\Lambda(d)\floor{\frac{x}{d}}\\
            &=\sum_{d\leq x}\Lambda(d)\left(\frac{x}{d}-\fract{\frac{x}{d}}\right)\\
            &=x\sum_{d\leq x}\frac{\Lambda(d)}{d}+\sum_{d\leq x}\Lambda(d)\fract{\frac{x}{d}}\\
            &=x\sum_{d\leq x}\frac{\Lambda(d)}{d}+\Oh\left(\sum_{d\leq x}\Lambda(d)\right)\\
            &=x\sum_{d\leq x}\frac{\Lambda(d)}{d}+\Oh(x).
        \end{align*}
        Since $S(x)=x\log x+\Oh(x)$ we get
        \[
            \log x+\Oh(x)=\sum_{d\leq x}\frac{\Lambda(d)}{d}+\Oh(x).
        \]
        \item[(2)]
        \begin{align*}
            \sum_{p\leq x}\frac{\log p}{p}&=\sum_{p^m\leq x}\frac{\log p}{p^m}-\sum_{\substack{p^m\leq x\\ m\leq2}}\frac{\log p}{p^m}\\
            &=\sum_{n\leq x}\frac{\Lambda(n)}{n}-\sum_{\substack{p^m\leq x\\ m\leq2}}\frac{\log p}{p^m}.
        \end{align*}
        Note that
        \begin{align*}
            -\sum_{\substack{p^m\leq x\\ m\leq2}}\frac{\log p}{p^m}&=\sum_{p\leq\sqrt{x}}\log p\sum_{2\leq m\leq\frac{\log x}{\log p}}\frac{1}{p^m}\\
            &\leq\sum_{p}\log p\sum_{m=1}^{\infty}\left(\frac{1}{p}\right)^m\\
            &=\sum_{p}\log p\left(\frac{p^{-2}}{1-p^{-1}}\right)\\
            &=\sum_{p}\frac{\log p}{p^2-p}<\infty\\
            &=\Oh(1).
        \end{align*}
    \end{enumerate}
}
\end{document}